% !Mode:: "TeX:UTF-8"

\chapter{延迟与容错：从Hystrix到Resilience4j}

\section{设计原理}

在微服务架构中，当某个服务出现故障或响应缓慢时，如果不加以控制，故障会通过调用链向上游传播，最终导致整个系统的"雪崩效应"。熔断器（Circuit Breaker）模式是预防雪崩效应的核心手段。

熔断器有三种状态：
\begin{description}
	\item[CLOSED（闭合状态）] 所有请求正常通过。当失败率超过阈值时切换到OPEN
	\item[OPEN（断开状态）] 请求直接失败（短路），执行降级逻辑。经过等待时间后切换到HALF\_OPEN
	\item[HALF\_OPEN（半开状态）] 允许部分请求试探通过。成功则回到CLOSED，失败则回到OPEN
\end{description}

Hystrix是Netflix开源的传统熔断器组件，但在Spring Cloud 2020.0.2中已被移除。本项目使用Resilience4j作为替代方案，应用于Gateway网关层。

\section{Gateway层熔断器配置}

\begin{lstlisting}[language=Java, caption=CircuitBreakerConfig.java]
@Configuration
public class CircuitBreakerConfig {
    @Bean
    public ReactiveResilience4JCircuitBreakerFactory
            reactiveResilience4JCircuitBreakerFactory() {
        ReactiveResilience4JCircuitBreakerFactory factory =
            new ReactiveResilience4JCircuitBreakerFactory();
        TimeLimiterConfig timeLimiterConfig = TimeLimiterConfig.custom()
                .timeoutDuration(Duration.ofSeconds(20))
                .cancelRunningFuture(true)
                .build();
        factory.configure(
            builder -> builder.timeLimiterConfig(timeLimiterConfig).build(),
            "globalCircuitBreaker", "userCircuitBreaker",
            "businessCircuitBreaker", "foodCircuitBreaker",
            "cartCircuitBreaker", "orderCircuitBreaker",
            "deliveryCircuitBreaker", "walletCircuitBreaker",
            "pointCircuitBreaker"
        );
        return factory;
    }
}
\end{lstlisting}

设计要点：
\begin{itemize}
	\item \textbf{独立熔断器}：为9个服务分别创建独立的熔断器，确保某个服务熔断不影响其他服务
	\item \textbf{统一超时}：所有熔断器配置20秒超时，超时后触发熔断
	\item \textbf{取消运行中的Future}：\verb|cancelRunningFuture(true)|确保超时后释放线程资源
\end{itemize}

\section{路由中的熔断器集成}

Gateway的每条路由规则都配置了CircuitBreaker过滤器和对应的降级处理地址：

\begin{lstlisting}[language=Java, caption=路由中的熔断器配置]
routes:
  - id: userServer
    uri: lb://elm-user-server
    filters:
      - name: CircuitBreaker
        args:
          name: userCircuitBreaker
          fallbackUri: forward:/userFallback
\end{lstlisting}

\section{降级响应设计}

FallbackController为每个微服务定义降级处理方法：

\begin{lstlisting}[language=Java, caption=FallbackController降级响应]
@RequestMapping("/userFallback")
@ResponseStatus(HttpStatus.FORBIDDEN)
public CommonResult<Map<String, Object>> userFallback(
        ServerWebExchange exchange) {
    Map<String, Object> metadata = new LinkedHashMap<>();
    metadata.put("fallback", true);      // 标记为熔断响应
    metadata.put("service", "user");     // 标识故障服务
    metadata.put("retryable", true);     // 标记可重试
    metadata.put("path", exchange.getRequest().getURI().getRawPath());
    return new CommonResult<>(403,
        "Gateway触发了用户微服务的熔断降级方法.", metadata);
}
\end{lstlisting}

降级响应携带丰富元数据：\verb|fallback: true|让前端判断这是熔断响应而非业务错误；\verb|retryable: true|提示可重试；\verb|path|和\verb|method|记录原始请求信息便于排查。

\section{容错机制分层总结}

系统的容错机制分为两层：
\begin{enumerate}
	\item \textbf{Gateway层熔断（Resilience4j）}：超时20秒触发，返回FallbackController降级响应，HTTP状态码403
	\item \textbf{Feign层降级（Fallback实现类）}：当目标微服务不可达时，调用Fallback方法返回兜底数据
\end{enumerate}

两层容错互为补充：Gateway层保障外部请求的可用性，Feign层保障内部服务间调用的鲁棒性。
